\subsection{The Local Distance Network}\label{sec:dn_definition}

The goal of measuring accurate and precise distances well into the Hubble flow 
($D \gtrsim 100 $ Mpc) to directly and empirically determine {\Hcst} requires the overlapping use of multiple techniques, a combination traditionally referred to as a ``distance ladder.''  The primary need for this approach is that geometric distance measurements have been limited in range ($D \lesssim 10$~Mpc, and most often $D \lesssim 100$~kpc), while long-range indicators, such as type-Ia supernovae (SNe~Ia), are too rare to provide enough examples within reach of geometric distances. 
Historically, measuring {\Hcst} thus required a rather large number of methods and combinations of inhomogeneous data sets \citep[e.g.,][]{2001ApJ...553...47F,2006ApJ...653..843S}. Significant improvements in precision and accuracy have since been achieved by streamlining distance ladder setups, focusing on high-quality and maximally homogeneous data sets while maintaining tight control of systematic errors by using the same instrument across rungs to cancel flux calibration offsets.
An example is the three-rung ``SH0ES'' distance ladder \citet{2009ApJ...699..539R}, which calibrates classical Cepheids (henceforth: Cepheids) as standard candles using distances measured by geometrical methods (i.e., higher precision MW parallaxes, detached eclipsing binaries in the Magellanic Clouds,the megamaser distance to NGC 4258), and, in turn, calibrates SNe~Ia luminosity using distances to their host galaxies measured using Cepheids observed exclusively using the {Hubble} Space Telescope (HST). Finally, SNe~Ia luminosity distances in the Hubble flow and their redshifts determine {\Hcst} \citep[see Fig.~10 in][]{2016ApJ...826...56R}. Such streamlined approaches provided significant benefits in terms of robustness and precision, in part because they avoided potentially correlated systematics, ultimately leading to the Hubble tension as discussed today.

In recognition that the tension may be indicating something profound, greater reliability has been sought through an increasing number of different methods, sources, and measurements. These provide multiple, interrelated constraints for the same or different astrophysical sources, resulting in partially independent paths to {\Hcst}. Many of these approaches replace certain steps with alternative methods, sources, or calibrators. For this reason, we believe it is appropriate to consider these tools in aggregate to comprise a ``{\DN},'' illustrated schematically in Fig.~\ref{fig:Metro}, to better convey the interdependence of these methods. While such a goal was considered ``lofty'' and potentially unreachable more than a decade ago \citep[cf.][and Fig.\,1 therein, originally credited to \citealt{2006pnbm.conf...79C}]{2013IAUS..289..351D}, the improvements to the systematics of several distance measurement techniques\hbox{---}inspired not least by the Hubble tension\hbox{---}now provide a wealth of robust information that allow to again ``diversify'' the base on which the local measurement of \Hcst{} rests.

The {\DN} provides two critical advantages on the path to a more accurate measurement of {\Hcst}: robustness (to reduce systematic errors) and statistical advantage (to reduce statistical uncertainties). Systematic errors can be recognized by the redundancy of methods allowing for analyses that ``leave one out.''  At the same time, redundancy offers the means to reduce statistical fluctuations through covariance-weighted averaging.  Informally, method combination and robustness has been evaluated through the display of ``whisker diagrams" which separate measurements of {\Hcst} by the combination of techniques they employ.  Intermediate measure comparisons have also been made, most importantly by comparing multiple ways to measure distances to specific SN Ia hosts, each calibrated by the same geometric source (e.g., Cepheids, TRGB, Miras, and JAGB), though these involve only a fraction of the available data \citep{2024ApJ...977..120R}. Formal covariance-weighting and combining has been attempted for only a limited set of indicators \citep{2022ApJ...934L...7R}. 

In this paper, we introduce an approach to the {\DN} that obviates these shortcomings.  By combining existing measurements \textit{at the component level}, rather than in terms of the resulting value of \Hcst, into a common, statistically rigorous framework encompassing a broad range of methods,  this approach yields a combined value of {\Hcst} with an uncertainty that reflects all available information.  Within this framework, we will also be able to include or exclude different subsets of measurements, thus identifying possible outliers. We will be able to inspect residuals at different levels, verifying whether they are consistent with their stated accuracies.  


\subsubsection{Nomenclature and definitions\label{sec:nomenclature}}

Given the intricate interrelations between different methods and measurements, and the complexity of the resulting framework, we define at the outset a set of terms, following past usage as closely as possible, that we will use to formulate our approach. The principles underlying the various methods alongside the datasets used are presented in Sect.\,\ref{sec:data_summary} and App.\,\ref{sec:data_appendix}.

\begin {description}
\item[Anchor] Any object, or collection of objects, whose distance is directly determined by geometric means, such as parallax or measurements of orbiting systems, and is used to calibrate the distances of other indicators. Anchors set the absolute scale of the \DN{}; all other distances, with a few exceptions (Section~\ref{sec:data_summary}), are measured {\it relative} to this scale.  Anchors used in this analysis include NGC 4258 (through Keplerian motion of circumnuclear masers), the Magellanic Clouds (through detached eclipsing binaries, DEBs), and the collection of Milky Way Cepheids (through trigonometric parallaxes).  The Milky Way uniquely provides distances to individual objects rather than a single extragalactic system, mostly measured by the ESA {\it Gaia} mission, which are then combined into a single calibration of the Leavitt Law for Galactic Cepheids.  To a lesser extent, depth effects are also present in the Large Magellanic Cloud (LMC) and Small Magellanic Cloud (SMC), but may be corrected through the use of empirical geometric modeling fit to the collection of DEBs. 
In the traditional distance ladder, these anchors of geometric measures are often referred as the first rung.

\item [Primary Distance Indicator]  An astronomical feature that can be measured or calibrated {\it directly} using the aforementioned geometric means.  Examples include the luminosity of the intercepts of the Leavitt Law of Cepheids or oxygen-rich Mira variables, as well as the luminosities of the Tip of the Red Giant Branch (TRGB) or the J-region of the Asymptotic Giant Branch (JAGB).

\item[Host] An object, typically a galaxy, whose distance can be estimated from its properties via a Primary Distance Indicator.  Relevant hosts include one or more Secondary Distance Indicators and by means of their identical distance, enable absolute calibration of the secondary indicator (i.e., converting relative distance to true distance).  In the distance ladder, these are often referred to as the {second rung}.

\item [Secondary Distance Indicator] An astronomical feature that can be measured in more distant systems (``hosts'') and ideally reach out to Hubble flow systems.  Secondary distance indicators used here include the luminosity of Type Ia and Type II supernovae (SNe~Ia, SNe~II), based on the measurement of objects in nearby hosts; the Tully-Fisher (TF) relation, which relates a spiral galaxy's luminosity to its velocity width; the standardized luminosity of Surface Brightness Fluctuations (SBF); and the Fundamental Plane (FP) of elliptical galaxies, calibrated using its properties in the Coma cluster. 

\item [Calibrator] An astronomical object used in the calibration of a Secondary Distance Indicator, such as SNe~Ia, SNe~II, and galaxies with luminosity estimated via TF or SBF relations.  A Calibrator is in a host (or \textit{is} a host, when the Secondary Distance Indicator is based on the whole galaxy), and its distance is constrained by the host distance.

\item [Group] A grouping (group or cluster) of galaxies that are assumed to be at a common distance, with appropriate dispersion due to depth effects.  Group membership is used to obtain a distance estimate for some Calibrators, including SBF calibrators in Virgo and Fornax and FP calibrators in Coma.

\item [Direct Distance Determination] A method that yields the distance to an astronomical object directly, without relying on intermediate calibration steps.  Such methods are used to determine the distance of anchors, via trigonometric parallaxes (in the Milky Way), the orbital signatures of red giant stars in DEB (LMC, SMC), or the relation between line-of-sight and angular velocity fields in a system of masers (NGC~4258). However, direct distance determinations can sometimes be applied to astronomical objects in the Hubble Flow (see next item) and thus provide a direct constraint on {\Hcst}. Examples include maser systems similar to that in NGC 4258, which yield an angular diameter distance by modeling their recession and angular velocity fields, and SNe~II calibrated via the Expanding Photosphere Method (EPM), which determines distance by comparing angular size measured from surface brightness and color temperature to physical size determined from the evolution of the velocity profile. 

\item [Hubble Flow System (also called Tracer)] An object at sufficiently large distance such that its cosmological redshift can be determined with good accuracy from its measured velocity; together with an angular or luminosity distance determination, such objects provide constraints on the value of $ H_0 $. Tracers used here include SNe~Ia, SNe~II, galaxies with luminosity distances estimated from TF or SBF relations, megamasers, and elliptical galaxies via FP.  The angular or luminosity distance can be determined using a Primary or Secondary Distance Indicator, or through a direct distance determination. Consistent with most determinations of the local value of {\Hcst}, we limit our analysis to redshifts large enough to reduce the impact of correlated flows due to large-scale structure, generally $z>0.01$ or $z>0.023$ and small enough ($ z \lesssim 0.15 $) that a simple, kinematic form of the redshift-distance relation suffices.  In practice, the effective redshift range for Hubble flow systems follows the relevant literature and depends on the tracer; tracers that require resolved galaxy images, such as SBF, typically occupy a lower redshift range than SNe~Ia.

\item [Use of Variance and Covariance] For each measured quantity, the original sources often define an uncertainty, which can be the combination (in quadrature) of several terms.  For example, the uncertainty in a TRGB-based measurement of the distance to a host galaxy (calibrator) may combine at least three terms: the uncertainty in the geometric distance to the anchor, such as NGC~4258; the uncertainty in the apparent magnitude of the TRGB in the anchor; and the uncertainty in the TRGB of the calibrator itself. However, several of these terms may be in common with other data. For example, all distance estimates anchored to NGC 4258 share the variance associated with its geometric distance measurements, and all TRGB distances estimated relative to NGC~4258 share the uncertainty of the apparent TRGB magnitude determination in this galaxy.
In order to make use of more than one such measurement, these terms must be included as \textit{covariances} between the relevant data, and thus added as off-diagonal elements to a full covariance matrix of the data system.  More details are given in the description of the equations in Section~\ref{app:equations}.

Based on the need to define covariance and successfully leverage a broad range of distance indicators, it is necessary to restrict the use of measurements to those with {\it direct} traceability to well-defined sources.  Likewise, direct linkages assume consistent measurements between sources (e.g., Cepheids in an anchor and an SN Ia host), a reasonable assumption when both employ the same telescope and instrument.  In contrast, the combination of measures from different telescopes and instruments involves substantial covariance (most certainly from different zeropoints) whose characterization is rarely provided in publications and is beyond the scope of this work to define.  The above serve as ``quality cuts'' for the inclusion of data in the distance network.  As an example, we make use of TRGB measurements in hosts calibrated with NGC 4258, where both are identically measured with HST (or JWST). However, we do not combine ground-based and space-based calibrations.


\end{description}

\begin{figure*}
    \centering
    \includegraphics[width=\linewidth]{Figures/LocalDistanceNetwork_Everything.png}
    \caption{The complete {\DN}, with all possible pathways illustrated. Anchors are objects that establish an absolute scale based on the methods shown to their left. Primary distance indicators (Cepheids, TRGB, Miras, JAGB) transfer the absolute scale to hosts (i.e., galaxies), the ensemble of which calibrates secondary distance indicators in the Hubble flow (tracers).  Exceptions are Megamasers and astrophysically modeled SNe~II, both of which serve as primary distance indicators and are capable of reaching the Hubble flow without intermediate steps. Green arrows illustrate direct connections between anchors or tracers and the method used to determine the absolute scale. Blue, violet, yellow, and red arrows show which calibrators constrain host distances; line width qualitatively distinguishes the attainable precision. Among hosts, rectangles qualitatively indicate overlap among objects measured via multiple methods.
    Diamond shapes represent groups. Dark gray arrows tie subsets of hosts whose distance is constrained by different calibrators to tracers. Any given arrow may represent multiple data sets, e.g., HST or JWST photometry of Cepheids, or optical vs.\ infrared photometry of SNe~Ia. The number of hosts is labeled for Cepheids, TRGB, JAGB, and Miras, with the number of hosts exclusively available to each method shown in parentheses.}
    \label{fig:DN_full}
\end{figure*}

\subsubsection {Distance Network Architecture}

Given very few assumptions,\footnote{We make conventional assumptions common to distance ladders: that the cosmological principle holds so that we do not live in a special place (i.e., there is no intrinsic difference in objects solely based on their separation from us); that evolution over short lookback times, of order a few hundred~Myr, is negligible; and that the expansion history for $z<0.15$ can be approximated with an expansion of a few kinematic terms, namely {\Hcst}\,, $q_0$\,, and $j_0$\,.} the measures and concepts given above may be linked together as a system of equations, as we will show below and in Appendix~\ref{app:equations}, which we refer to as the {\DN}. The multiplicity of data (see Sect.~\ref{sec:data_summary} and Appendix~\ref{sec:data_appendix}) means we can optimize the system by introducing free global parameters.   The wealth of sources, methodologies, and distance indicators also provide important additional and complementary  information which can help strengthen the distance ladder or can contribute to constrain {\Hcst}.  For example, the distance to a given supernova host can be measured via different methods (TRGB, Cepheids, JAGB, Miras) depending on which observations are available for that galaxy, and each method can be related to a different anchor.  However, these determinations are not independent, as explained under ``Use of Variance and Covariance'' in Sect.\,\ref{sec:nomenclature}. The full {\DN} is illustrated in Fig.~\ref{fig:DN_full}.

Within the approximations of the present analysis, all equations are linear (or linearized), and all probability distributions are Gaussian in magnitude, distance modulus, or $\log_{10}{(H_0)}$.
Due to shared calibration sources among primary or secondary distance indicators, a covariance matrix with off-diagonal elements provides a useful method for proper accounting and weighting. The optimization procedure is that of a generalized linear least squares problem. The global solution to this problem provides best-fit values for the underlying parameters of the {\DN} as well as their uncertainties, most notably the distance moduli of included host galaxies, absolute magnitude calibrations, and the Hubble constant. These parameters include $\log_{10}{(H_0)}$, the distance moduli to host systems, $\mu_{\mathrm{host},\mathcal{H}}$, and groups included in the host data, $\mu_{\mathcal{G}}$, as well as the reference magnitudes for each type of calibrator, $M_{\mathrm{ref},\mathcal{T}}$. % \nilsc{I don't know what is the added value by showing symbols here.}
The full description of the system of equations and the covariance matrix is given in Appendix~\ref{app:equations}. The employed data sets are listed in Table~\ref{tb:datasets}, described briefly in Section~\ref{sec:data_summary}, and in greater detail in Appendix~\ref{sec:data_appendix}.  

\begin{deluxetable*}{lll}[t]
\tablewidth{0pc}
\tabletypesize{\scriptsize}
\tablecaption{Distance Network Datasets. Note that for all references we have made sure to decompose the listed uncertainties into those coming from the anchor (which the {\DN} treats separately) and those coming from the actual measurement of the distance indicator (separated to that within and without the anchor).
\label{tb:datasets}}
\tablehead{\colhead{Measure} & \colhead{Reference} & \colhead{Notes}}
\startdata
\hline
Milky Way (MW)& \citet{2018ApJ...855..136R, 2021ApJ...908L...6R,2022ApJ...938...36R} & (a)\\
LMC& \citet{2019Natur.567..200P} &$\mu_0$=18.477$\pm$0.024; converted from distance in kpc\\
SMC& \citet{2020ApJ...904...13G} & $\mu_0$=18.977$\pm$0.032; converted from distance in kpc\\
NGC 4258 Masers& \citet{2019ApJ...886L..27R} & $\mu_0$=29.397$\pm$0.032\\
\hline
\multicolumn{3}{c}{Primary distance indicators} \\[-0.015cm]
\hline
HST Cepheids & \citet{2022ApJ...934L...7R} & SN hosts, separate measures, N4258, MW, LMC (R22a fits 9,10,11) \\
& \citet{2022ApJ...938...36R} & MW cluster Cepheid photometry \\
& \citet{2019ApJ...876...85R} &  LMC  \\
 & \citet{2022ApJ...940...64Y} &  NGC 4258  \\
  & \citet{2024ApJ...973...30B}  &  SMC  \\
JWST Cepheids & \citet{2024ApJ...977..120R,2025arXiv250901667R}  & \\
JWST TRGB & \citet{2024ApJ...966...89A}   & \\
& \citet{2024ApJ...976..177L, 2025arXiv250408921L}  & \\
JWST TRGB & \citet{2025arXiv250311769H}  & \\
JWST TRGB & \citet{2024ApJ...973...83A,2025ApJ...982...26A} &  \\
HST TRGB & \citet{2024ApJ...966...89A}  &  \\
HST TRGB & \citet{2025ApJ...985..203F} &  \\
HST TRGB & \citet{2022MNRAS.514.4620D}  & and contained references \\
JAGB & \citet{2025ApJ...988...97L} & \\
JAGB &  \citet{2025ApJ...985..203F}  & \\
Miras & \citet{2024ApJ...963...83H} & Modified in correspondence with C. Huang to N4258-only values \\
\hline
\multicolumn{3}{c}{Secondary distance indicators} \\[-0.015cm]
\hline
SNIa Pantheon+ & \citet{2022ApJ...938..110B}  & \url{https://github.com/PantheonPlusSH0ES/DataRelease}\\
SNIa SNooPy (post v2.7) & C. Burns 2025, priv.\ comm.  &  Re-fitting the official data from \citet{2024ApJ...970...72U} \\
SNIa SNooPy (pre v2.7) & \citet{2024ApJ...970...72U}  & \url{https://github.com/syeduddin/h0csp/tree/main/data/working} \\
SNIa BayesSN & \citet{2023MNRAS.524..235D} & also E.~Hayes 2025, priv.\ comm.\ for additional fits\\
SNIa Salt3 & \citet{2021ApJ...923..265K}  &  \\
SNIa IR & \citet{2023AA...679A..95G}  & J, H bands separately\\
SNII (SCM) & \citet{2020MNRAS.496.3402D}  & T. de Jaeger 2025, priv.\ comm.\\
TF & \citet{2024MNRAS.531...84B} & Cosmicflows-4 catalog; data curated by the authors\\
\hline
\multicolumn{3}{c}{Other} \\[-0.015cm]
\hline
SNII EPM & \citet{2025AA...702A..41V},  & Removed galaxies (M61 and N6946), where peculiar \\
 & \citet{2023AA...672A.129C} &  velocities could not be meaningfully established.\\
% &   &  & \\
Megamasers & \citet[Tab.~1]{2020ApJ...891L...1P} & We also adopt our own peculiar velocity treatments\\
Coma FP & \citet{2025MNRAS.539.3627S}  & See also Appendix \ref{sec:eq_special}\\
SBF & \citet{2025ApJ...987...87J}  & Ancillary info from priv.\ comm.\ with J.~B.~Jensen.\\
\enddata
\tablecomments{(a) For the measurement of calibrators within the Milky Way galaxy, its extent is quite relevant --- therefore we adopt the individual measurements standardized to a 1~kpc reference distance.} 
\end{deluxetable*}

\subsubsection{Selecting the baseline}\label{sec:define_baseline}

There was strong consensus among the workshop participants, which is also shared more broadly within the community, that not all data, measurements, and tools are viewed with equal confidence.  For this reason, we decided to produce a ``baseline'' result based on all methods which enjoy the broadest and highest level of confidence (sometimes referred to in the literature as ``gold standards''), along with variants to the baseline that incorporate additional methods or datasets, or exclude certain data in order to reflect a broad range of minority viewpoints and explore the robustness of the results. To reach a true consensus for the baseline and variants, we adopted a collaborative, methodologically inclusive approach during the ISSI workshop ``What’s under the H0od.'' Prior to analyzing any combined fits, we engaged in extensive discussion of the various available methods, the choice of data sets (for instance, Cepheids from HST, TRGB distances from JWST, etc.), distance calibrators, and the systematic uncertainties involved in the determination of local distances and \Hcst\,. These included both classical and emerging techniques, ranging from geometric anchors (e.g., megamasers, {\it Gaia} parallaxes, detached eclipsing binaries) to stellar standard candles (e.g., Cepheids, TRGB, Miras, JAGB), as well as other methodologies such as Surface Brightness Fluctuations from HST and JWST and the Fundamental Plane from DESI. We also carefully reviewed the use of Type~Ia and Type~II supernovae, which extend the distance ladder into the Hubble Flow. In particular, we examined various SNe~Ia data sets (e.g., Pantheon+, CSP, NIR-only samples) as well as SNe~II determinations based on both the standard candle method and spectral modeling.

\begin{deluxetable}{llcccc}[h]
\tablewidth{\textwidth}
\tablecaption{Result from anonymous vote on ``technical readiness level" of distance calibrators. Above the line the consensus is that these  will be included in the baseline, below the line will be included  as variations.}
\label{table:rank1}
\tablehead{\colhead{Rank} & \colhead{Distance calib}& \colhead{in baseline}&\colhead{in variants} & \colhead{excluded} & \colhead{abstain}}
\startdata
\hline
1&Cepheids &  {\bf 28} & 1 &0&0 \\
2&DEB     & {\bf 26}   &3  &0 &0\\
3& TRGB   &  {\bf 26}&  2   &0 &1\\
3& NGC4258 & {\bf 26}&   2&  1&0\\
4& {\it Gaia} parallaxes & {\bf 24}& 3&2&0\\
\hline
5& Miras & 8& {\bf 15}&1&3\\
6& JAGB & 5& {\bf 15}&3&3\\
\hline
\enddata
\end{deluxetable}

\begin{deluxetable}{llcccc}[h]
\tablewidth{\textwidth}
\tablecaption{Result from anonymous vote on the method and secondary distance indicators.  Above the line the consensus is that these  will be included in the baseline, below the line will be included  as variations.}
\label{table:rank2}
\tablehead{\colhead{Rank} & \colhead {Method} & \colhead{in baseline}& \colhead{in variants} & \colhead{excluded} & \colhead{abstain}}
\startdata
\hline
1&SNIa &  {\bf 32} & 0 &0&0 \\
2&SBF     & {\bf 23}   &8  &0 &1\\
3& Masers in Hubble flow   &  {\bf 19}&  12   &0 &1\\
\hline
4& SNe~II, empirically standarized &15&   {\bf 14}&  0&3\\
5&SNe~II, astrophysically calibrated & 1& {\bf 22}&3&6\\
6& DESI FP (calibrated to Coma) & 11& {\bf 20}&0&1\\
7& TF from CF4 & 8& {\bf 21}&1&2\\
\hline
\enddata
\end{deluxetable}

A key element of this process was a series of open, expert-led discussions in a plenary setting, during which participants evaluated the strengths, systematics, data availability, and limitations of each method and dataset. Following these in-depth assessments, we held anonymous ballot votes to determine the configuration of the \textbf{baseline} analysis and well-motivated \textbf{variants} thereof, as well as configurations that should be excluded entirely.

The first vote identified the primary and secondary distance indicators and methods considered most technically mature and broadly supported, cf.  Table~\ref{table:rank1}. In a second vote, we decided on specific combinations, methodological choices, and alternative calibrations, resulting in a comprehensive suite of variants, see Table~\ref{table:rank2}. Methods were included in the baseline if they garnered more than half the total votes, including abstentions, as indicated by the horizontal lines in Tables~\ref{table:rank1} and \ref{table:rank2}.  All votes were taken prior to deriving a value of {\Hcst}, and their outcomes were considered binding. The datasets included in the baseline configuration are listed in Table~\ref{tb:datasets} and discussed in detail in Appendix~\ref{sec:data_appendix}. Where multiple datasets were available in principle, preference was given to datasets that had been fully published 
and/or have been most widely adopted. Alternative datasets, notably of SNe~Ia (i.e., Pantheon$+$ vs. CSP vs. BayesSN), were included as variants.

This process aimed at identifying the most widely supported and robust measurement paths based on current data and understanding, not at enforcing a single value of \Hcst. The resulting baseline, variants, datasets and code provide a flexible framework to assess the stability and reliability of the derived consensus {\Hcst} value under well-motivated methodological changes. This structured approach reflects the collective judgment of the experts present at the workshop, and provides a reproducible path forward for future analyses. Hence, ``consensus'' here refers to an a priori agreement on \emph{how} a {\DN} should be constructed in terms of methods, datasets, and their combinations, rather than on the outcome, which was not known when the process was fixed.

In summary, the baseline analysis reflects the scientific consensus among workshop participants concerning the optimal combination of datasets and methodologies as established prior to obtaining results. Variants of the baseline analysis were performed to test the robustness of the result to individual methods, to different datasets corresponding to identical methodologies, and to quantitatively explore ``hypothesis-driven'' setups that combine multiple elements deviating from the baseline. Further detail on the process of selecting variants is provided in Sect.~\ref{sec:select_variants} below.

\subsubsection{Selecting variants and other methodological choices }\label{sec:select_variants}

We defined analysis \emph{variants} to explore the sensitivity of the {\Hcst} determination to {\it well-motivated} modifications in the analysis setup. Possible variants were identified through brainstorming and subsequently discussed for their merits and weaknesses. These variants serve as internal consistency tests (or ``null tests''), helping to assess the robustness of the baseline result under plausible changes to inputs or methodology. Considerable discussion was involved in ensuring that variants were well-motivated prior to knowing their impact on \Hcst{}. We specifically avoided combinations targeting maximal changes in \Hcst{}, as these would require {\it a posteriori} information or intermediate outcomes, potentially becoming prone to confirmation bias. Our process described in Sections~\ref{sec:philosophy} and \ref{sec:define_baseline} was key to taking these decisions well informed by broad interdisciplinary and expert knowledge. In each variant, we furthermore aimed to adequately reflect the current state-of-the-art in the various disciplines by including all data sets sufficiently precise to carry significant weight in the determination of \Hcst{}. In other words, we omitted data sets or methods only when motivated by a specific hypothesis in order to avoid artificially inflated uncertainties (i.e., caused merely by insufficient data) that would reduce the Hubble tension in an uninformative way not representative of the state of knowledge.

We reiterate that most variants are effectively ``null tests.''  Their consistency (or lack thereof) serve as a test of robustness of the baseline.  In general, they are not fair substitutes for the baseline solution, and should always be referenced together with a description of their underlying choices.  It is also critical to note that the values of {\Hcst} obtained with all variants are \textit{highly correlated}.  They share a large fraction of the underlying data, and are therefore \textit{expected} to differ from each other by much less than the nominal uncertainty.  Therefore, quantifying if individual datasets have anomalous pulls on the final result requires a more careful analysis; see Section~\ref{sec:variants} for further discussion. 

The variants discussed in this paper can be organized into different categories:

\paragraph{Add-one-in variants} include additional information that was not voted for inclusion in the baseline, introduced through a step-by-step process. This category can  introduce additional anchors (i.e., the SMC), as well as primary or secondary distance indicators. For example, V01 (Baseline+JAGB) adds JWST/NIRCam distances of the JAGB in supernova hosts calibrated to that in NGC~4258. The addition of calibrators to the baseline configuration serves to explore whether these supplementary methods shift or reinforce the baseline results. Furthermore, this category includes the incorporation of alternative distance indicators in the Hubble flow, such as SNe~II calibrated via the standard candle method or spectral modeling, the Fundamental Plane from DESI anchored to Coma, and the Tully-Fisher relation as implemented in the Cosmicflows-4 catalog \citep[CF4;][]{2020ApJ...902..145K}. 

\paragraph{Leave-one-out variants} involve the exclusion of individual methods or calibrators from the baseline configuration, motivated by the hypothesis that there is an undiscovered error in a type of measurement. Each variant is constructed by removing one key element at a time—such as Cepheids, SNe~Ia, TRGB, {\it Gaia} parallaxes, or NGC~4258—and observing the resulting impact on the derived {\Hcst} value. These tests help identify whether any individual dataset or calibrator is significantly impacting the result. This category also includes the exclusion of specific indicators in the third rung, such as SBF, SNe~Ia, and masers in the Hubble flow.  Since leaving out a given element can implicitly remove other elements (e.g., specific SNe~Ia calibrators), we have constructed custom baselines (cf.~below) to separately assess the effect of the implicit (undesired) modification relative to the baseline.

\paragraph{Hypothesis-driven variants} involve compound 
configurations motivated by specific physical or observational hypotheses. These include the use of alternative calibrations (e.g., CSP/SNooPy instead of Pantheon+/SALT2 for SNe~Ia), different treatments of systematics (e.g., inclusion or exclusion of peculiar velocity corrections, metallicity corrections, or near-infrared SN data), and restricted subsets of the data (e.g., only modern SNe~Ia, or cutting the Hubble flow sample at $z > 0.06$). Additionally, we included a variant that considered the Cepheid Leavitt law to be independent of chemical composition.

\paragraph{Instrument-suspicious variants} explore the impact of removing specific observatories. In particular, we tested the exclusion of all HST or JWST-based observations, all {\it Gaia} measurements, etc.

\paragraph{Custom baselines} The exclusion of certain measurements can also cause the exclusion of some SNe~Ia calibrators; for example, when Cepheid measurements are excluded, only 35 of the 55 SNe~Ia calibrators in the baseline can be included, since the other 20 only have determinations of the host distances through Cepheids. The resulting change in the {\Hcst} value and uncertainty is due to both the exclusion of the measurements and the change in the SNe~Ia calibrator sample. To cleanly separate the two effects, we define a custom baseline that includes those measurements, but is restricted to only the SNe~Ia calibrators for which other measurements are available. Custom baseline versions are not variants {\it per se}; they are only intended to facilitate the interpretation of the results for the corresponding variants.  

\paragraph {Include Everything}
We also consider a variant in which all independent methods and data sets are included, to illustrate the accuracy that can be achieved with presently available measurements if systematic effects and other issues are resolved.  A secondary variant in this class uses all independent methods \textit{except} Tully-Fisher, for which the current data set has excess dispersion.  In either version, this variant includes only one set of measurements for SNe~Ia, since different measurements are likely correlated (because of astrophysical variance) to a degree that has not been sufficiently quantified in the literature.  

\paragraph {Additional solutions}  In addition to the variant solutions described above, we also consider special solutions designed to test our methodology. These are not ``variants'' in the same sense as those listed previously, and do not appear in Table~\ref{tab:variants}; they are discussed in mode detail in Sections~\ref{sec:orthogonal_paths}, \ref{sec:consistency_checks}, and Appendix~\ref{sec:comparisons}.  They include \textit{independent grouping solutions}, \textit{consistency checks}, and \textit {emulator solutions}.  Independent grouping solutions consist of configurations in which two or more fully independent paths through the distance ladder are identified, sharing no common calibrators or intermediate steps. These allow for the construction of entirely separate determinations of {\Hcst}, such that agreement between them provides a strong consistency check and minimizes the risk of shared systematic effects. They are the analogous to splitting the data in uncorrelated halves or thirds. Two such variants are discussed in Section~\ref{sec:orthogonal_paths}.  Consistency checks are employed to make sure that the network results are statistically self-consistent for different paths through it and are further discussed in Section~\ref{sec:consistency_checks}.  Emulator solutions involve a set of \textit{emulators} designed to reproduce key published results, such as SH0ES, CCHP, and recent SBF analyses, within our unified framework and dataset handling. These configurations serve both as validation tests of our pipeline and as transparent benchmarks for comparison with previous literature, and they may also help explain sources of differences. These are discussed in Appendix~\ref{sec:comparisons}.

